\title{Organizing Intelligence}
\begin{document}
\maketitle

\begin{abstract}
When will advances in AI lead to economic transformation? We develop a model of how intelligence, the ability to achieve goals, matters for organizations, the key unit of economic production. Agents (human or AI) with multidimen- sional capabilities sort endogenously into competing firms. A firm’s productivity depends on its agents’ capabilities; we distinguish competence, which shifts pro- ductivity levels, from coordination, which governs how productivity scales with size. Two differences between AI and humans drive transformation: AI expertise is reproducible while humans’ are in fixed supply, and AIs may not face the same limits for coordination as humans. We analyze how equilibrium market structure, productivity, employment, and wages evolve as AI capabilities improve. When AI is more capable than the average human, organizations of complementary humans can nonetheless outperform AI. Increases in AI competence can change the composition of workers within firms and decrease firm size. As AI coordi- nation improves, however, this can break the Coasean ceiling on firm size and the market structure undergoes a phase transition toward large, concentrated AI-heavy firms. Changes in human employment are gradual when AI displaces workers within mixed human-AI firms, or sudden when AI firms outcompete human firms. We decompose wage changes into a technology effect, which changes the value of humans within existing firms, and a reorganization effect as firms adjust size and compos
\end{abstract}

\section{Recovered Text}
ORGANIZING INTELLIGENCE*
Alex Imas†
Andrew Koh†
Shane Legg
Google DeepMind
October 29, 2026
Abstract
When will advances in AI lead to economic transformation? We develop a
model of how intelligence, the ability to achieve goals, matters for organizations,
the key unit of economic production. Agents (human or AI) with multidimen-
sional capabilities sort endogenously into competing firms. A firm’s productivity
depends on its agents’ capabilities; we distinguish competence, which shifts pro-
ductivity levels, from coordination, which governs how productivity scales with
size. Two differences between AI and humans drive transformation: AI expertise
is reproducible while humans’ are in fixed supply, and AIs may not face the same
limits for coordination as humans. We analyze how equilibrium market structure,
productivity, employment, and wages evolve as AI capabilities improve. When
AI is more capable than the average human, organizations of complementary
humans can nonetheless outperform AI. Increases in AI competence can change
the composition of workers within firms and decrease firm size. As AI coordi-
nation improves, however, this can break the Coasean ceiling on firm size and
the market structure undergoes a phase transition toward large, concentrated
AI-heavy firms. Changes in human employment are gradual when AI displaces
workers within mixed human-AI firms, or sudden when AI firms outcompete
human firms. We decompose wage changes into a technology effect, which changes
the value of humans within existing firms, and a reorganization effect as firms adjust
size and composition. Reorganization can reverse the technology effect on wages.
*We thank Allan Dafoe, Matt Elliott, Joshua Gans, Gillian Hadfield, Andrea Prat, Eduard Talamas,
Phil Trammell, and Birger Wernerfelt for helpful conversations and suggestions. We also thank Pedro
Aldighieri for superb comments and research assistance. Various AI systems were used for assisting
with mathematical arguments as well as proofreading for consistency and clarity. The paper expresses
the views of the individual authors and does not necessarily reflect those of their affiliated institutions.
This research is intended solely for academic discussion and does not represent corporate forecasts,
product roadmaps, or empirical statements regarding Google’s current or future operations.
†Lead authors with equal contribution listed in alphabetical order; emails: imasa@google.com and
andrewkoh@google.com. Imas is also affiliated with The University of Chicago; Koh is also affiliated
with Columbia University.
1 INTRODUCTION
When will AI transform the economy? What kinds of firms will emerge, and what will
happen to the returns to human expertise and to wages? We develop a model of how
the capabilities of human and artificial intelligence interact with economic forces to
jointly determine the composition and size of firms, surplus, and human wages. We
model intelligence as the ‘ability to achieve goals’ (Legg and Hutter, 2007). This has
typically been measured through isolated benchmarks that evaluate a single agent
in isolation. But the economically relevant unit is the organization that competes in
equilibrium under the forces of supply and demand (Samuelson, 1948).
Economic models often treat new technology either as a tool that makes workers more
productive or as a factor that substitutes for humans task-by-task. But AI is different
in three conceptually important ways. First, human capabilities are fixed, which can
generate Ricardian scarcity rents in the form of higher wages for increasing expertise;
by contrast, AI capabilities can be copied at the cost of compute which eliminates these
rents. Second, human capabilities differ in jaggedness across people while copies of AI
models overlap in their jaggedness; sorting heterogeneous workers into organizations
creates what we call mosaic intelligence, such that human teams can out-compete
more competent AI. Finally, unlike machines of the past that had to operate within
human firms, AI agents can coordinate directly with one another to form standalone
organizations. This can lead to rich non-monotonicities in market structure as AI
capabilities increase.
We build a simple theoretical model incorporating those three features to formally
characterize potential economic scenarios.1 For instance, as AI models become more
competent at specific tasks, firms might first shrink as human managers can do more
with AI agents. But past a coordination threshold, AI can break the Coasean ceiling
and trigger a sudden phase transition toward a giant AI firm. Beyond this ceiling, firm
sizes start to shrink continuously once again as fewer AI agents suffice to produce
enough output to saturate the market.
In our model, each human and machine is endowed with a multidimensional capa-
bility profile. Multidimensionality captures an agent’s ability to achieve goals across
different domains.2 Such agents—both human and machine—sort into firms whose
output depends on both size, or how many workers the firm employs, and composi-
1As we take pains to emphasize throughout the paper, these scenarios are not empirical forecasts
but outcomes based on the model’s assumptions.
2See, e.g., Dell’Acqua et al. (2026); Gans (2026).
1
tion of the capabilities of its constituent members. A central distinction is between
competence, which captures domain-specific capabilities such as robotics, writing, or
logical reasoning, and coordination, which captures capabilities that allow agents to
work better with each other. This might comprise a better understanding of norms,
reduced communication costs, reduced agency frictions, and so on. A firm with more
competent agents enjoys a level shift in output. A firm with more cooperative agents,
on the other hand, enjoys better returns to scaling—it can function better as the orga-
nization grows in size. In equilibrium, human and AI agents endogenously sort into
competitive firms that serve a downstream market.3
How does the equilibrium market structure change as AI systems improve? Where are
humans employed and how much do they get paid? This hinges on the direction of
improvement, and whether specific capability thresholds are crossed. Our framework
helps rationalize the why the rapid increases in AI capabilities across a wide array
of domains can be accompanied by relatively little economic transformation. While
current AI models are beginning to match or exceed human performance across many
domains, they are still fairly poor at coordinating. Indeed, today’s labor market reflects
such capabilities: no big effects on employment or market structure, but leaner startups,
flatter org charts, and steady or rising demand for experienced managers who can
direct complex work across teams.4 By taking seriously the endogenous organization
of firms at different scales, our model makes nuanced and rich predictions on how
different directions of AI improvements will shape surplus, employment, and wages
and, in so doing, clarifies what it would take for economic transformation.
2 AN ILLUSTRATIVE EXAMPLE
Before turning to the general model, we build up a simple and concrete economy
step-by-step to illustrate the key mechanisms of the paper.
3Our equilibrium concept requires that every industry is competitive: producers and workers take
the offered prices and wages as given, active firms break even, no feasible entry is profitable, and
markets clear. Moreover, no entrant can profitably reorganize the industry under another competitive
continuation while making all workers and consumers weakly better off (via transfers). This will allow
us to make clear predictions on how the economy might transition between fragmented industries
(with infinitesimal firms), concentrated industries (with large firms), and mixtures thereof. In special
cases with ‘interior’ firms, our equilibrium concept coincides with competitive equilibrium.
4For superhuman capabilities see Jumper et al. (2021); Kwa et al. (2025); for muted economic effects
see Yotzov et al. (2026).
2
2.1 The
economy.
Each
worker
(human
or
AI)
brings
capability
vector
y = (y1, y2, y∗), combining two productive task skills—coding y1 and writing
y2—with a coordination capability y∗.
Let H denote the human population of
NH = 14 workers across three roles (Table 1). Coders and writers are mirror-image
specialists: each excels in one domain and struggles in the other. Coordinators are
weak at both tasks but excel at facilitating cooperation across people. Any person
who stays out of production takes an outside option worth w = 1/4. AI enters
in Section 2.3: firms can run any number of copies of an AI agent with profile
yM = (yM
1 , yM
2 , yM
∗) at a compute cost of r = 1 per copy.
Role
Task skills
Coordination
Supply
Coder
9, 1
3/4
6
Writer
1, 9
3/4
6
Coordinator
1, 1
0.97
2
Table 1: Capabilities of agents in the economy. Task skills report coding y1 first and
writing y2 .
Competence. Competence governs how much output a firm produces per worker at a
given size. A firm’s competence L is the harmonic mean of its members’ average skills
( ¯y1, ¯y2) (Kremer, 1993; Jones, 2011):5
L =
2
1/ ¯y1 + 1/ ¯y2
.
Figure 1(a) plots this aggregator in skill space. Each curved contour traces skill pairs
that deliver the same competence. Because the harmonic mean penalizes bottlenecks,
a firm is held back by its weaker domain: solo coders and solo writers sit near opposite
axes on a low contour (9/5 = 1.8) despite excelling at one task, while a coordinator
working alone sits on the lowest contour (1). Scored in isolation without complemen-
tary sorting across people, the human population averages standalone competence
and coordination
¯LH =
1
NH ∑
i∈H
L(yi
1, yi
2) ≈1.69
and
¯yH
∗=
1
NH ∑
i∈H
yi
∗≈0.78.
5The harmonic mean corresponds to a symmetric constant-elasticity-of-substitution aggregator
L =
 1
2 ¯yρ
1 + 1
2 ¯yρ
2
1/ρ with ρ = −1, giving an elasticity of substitution of 1/2. Following Kremer (1993)
on complementary O-ring production tasks and Jones (2011) on weak links in intermediate inputs, an
elasticity strictly below unity ensures that shortfalls in any essential domain bottleneck output without
introducing strict Leontief kinks; see also Aghion et al. (2019).
3
1
5
9
coding
1
5
9
writing
capability
envelope E(H)
L = 1
L = 9=5
L = 5
coder
writer
coordinator
coder–writer team
(a) Capability envelope vs. competence
1
6
16.5
50
100
firm size (workers, log scale)
0
5
10
15
20
output per worker
coders only
coder–writer
team
higher
coordination
supercooperative
(b) Scale and firm size
Figure 1: The capability envelope, competence, and firm scale.
Coordination and firm size. Coordination governs how productivity changes as a firm
adds workers. As Coase famously asked, why doesn’t one giant firm employ everyone?
Larger firms divide work more finely, lifting output per worker, yet each person must
also coordinate across more colleagues. Bounded human cognition and scarce attention
mean that information degrades across handoffs and key details go unheeded (Simon,
1947, 1971; Marschak and Radner, 1972); incentives across individuals also diverge as
group size grows (Williamson, 1967; Holmstr¨om, 1982; Becker and Murphy, 1992).
Consider a team building and documenting a software product. Two people working
side by side can hold the full design in mind and resolve mismatches immediately. But
expand that team to twenty people split across subteams, and members spend more
time in meetings, meaning gets lost in handoffs, and subteams drift toward individual
goals rather than the team’s mission, i.e., misaligned incentives. Classic span-of-
control theory formalizes how finite managerial attention bounds firm size through
these limits to coordination (Kaldor, 1934; Williamson, 1967; Calvo and Wellisz, 1978;
Lucas, 1978; Qian, 1994). Garicano (2000) shows how hierarchies solve this issue by
economizing on scarce attention by routing tasks within the firm.
Output per worker F/m in a firm of m workers combines these forces:
F
m =
L
|\{z\}
Firm
comp.
× eϕm/(1+m)
|
\{z
\}
dividing work
× (1 + m) ¯y∗
|
\{z
\}
building on
colleagues
×
1
1 + m
| \{z \}
split
attention
where ϕ = 7
4.
4
Dividing work raises productivity quickly at small scale and then levels off. Mean-
while, each worker builds on colleagues’ output at rate (1 + m) ¯y∗, governed by the
firm’s average coordination ¯y∗, while attention is strained in proportion to 1/(1 + m).
Of course this specification is a little ‘ad hoc’ but it will make the numbers concrete;
our main model will impose more general assumptions on how the capabilities of
individual workers matter for firm scaling.
Whenever ¯y∗< 1, bounded attention and potential misalignment eventually outweigh
what each additional worker contributes to colleagues. At small scale, the steep gains
from dividing work dominate that net coordination friction, so output per worker
rises initially. Once division-of-labor gains level off, coordination friction takes over.
Given these forces, output per worker peaks at the firm’s best size
m∗( ¯y∗) =
ϕ
1 −¯y∗
−1,
and declines thereafter (Figure 1(b)). The key takeaway here is that higher competence
shifts the level of the curve up while leaving best size unchanged, whereas better
coordination increases returns to scaling, shifting best size to the right—the optimal
size of the firm increases with coordination.
When ¯y∗> 1, the gains from building on colleagues’ work exceed the losses from
splitting attention. Output per worker rises continuously with firm size rather than
peaking and turning down (the dashed curve in Figure 1(b)); a larger firm always
produces more per worker than a smaller one. For instance, Radner (1993) shows that
organizing information processing efficiently within the firm can generate increasing
returns to scale. We call capability above this threshold supercooperative.
Human organizations run into decreasing returns for at least three reasons. First,
human attention and memory are finite, creating constraints on the ability of agents to
coordinate (Simon, 1971; Brooks, 1975). Second, much of human knowledge is tacit
(Polanyi, 1966), passing work from one person to another compresses rich internal
context into lossy summaries (Graeber et al., 2026). Third, individual humans have
private career concerns with their own objective functions, so incentive misalignment
can compound with scale (Holmstr¨om, 1982).
Reproducible AI agents need not share these constraints and frictions. Copies of a
single model share identical weights and representations, can pass complete working
state—entire codebases, intermediate derivations, and full context windows—without
lossy translation, and execute a common objective. In large-scale software engineering
or automated scientific discovery, thousands of parallel instances can explore distinct
5
branches simultaneously while sharing verified state in real time. Whether and when
AI crosses ¯y∗= 1 is an empirical question (see Ord, 2026, for some early experiments).
Markets and equilibrium. Firms sell output into competitive markets with inverse
demand
P(Q) = 0.3

1 −
Q
200N

,
where Q is total market output and N measures market size.6 Our economy comprises
two markets: a deep market (N →∞), where Q/(200N) →0 so buyers absorb any
finite quantity at price 0.3, and a niche market (N = 1), where buyer willingness to pay
falls to zero at Q = 200 so finite market size can cap how large a firm grows. In our
equilibrium concept (to be specified later) active firms break even, no entrant can earn
positive profit, and competition selects the industry structure (collection of firms) that
maximizes total surplus.7 Each worker’s equilibrium wage equals the shadow value
of how much she contributes to total surplus.
2.2 A sequence of human-only economies. Before we introduce AI, we will showcase
our model through a sequence of increasingly complex human economies.
One human type. We begin with one human type with just six coders. Humans can
work in either market, so all output sells into the deep market where demand holds
price at 0.3, leaving the niche market unserved for now. Weak writing holds a coder’s
competence low at 9/5 = 1.8 (Figure 1(a)), while human coordination (3/4) places
best firm size at 6 (the bottom curve in Figure 1(b)). Grouping all six coders into a
single firm lets them divide work and produce 5.0 per person (29.8 in total), compared
with 3.6 per person (21.8 in total) working alone as solo contractors. Yet coordination
strictly caps firm size: if twice as many coders entered, in equilibrium they would
split into two separate firms of six producing 59.5 in total (5.0 per coder) rather than
merging into one twelve-person firm producing 57.2 in total (4.8 per coder), where
coordination friction outweighs division of labor, as shown along the coders-only
curve in Figure 1(b).
Because consumers in the deep market want far more output than six coders can
supply, each coder commands a substantial wage premium—1.49, or roughly six times
her outside option (Figure 2(b)). The basic takeaway here is that coordination caps
6Units of output are normalized so that running one AI copy costs r = 1; wages and outside options
are measured in the same units.
7We will allow for non-integer firms to convexify the economy which guarantees the existence of an
equilibrium; see the main model for a discussion.
6
how large a firm can grow and that fixed human supply turns scarce talent into a
substantial wage premium over the outside option.
Two human types: mosaic intelligence. Now add the six writers, whose skills mirror
coders in reverse. Each coder and each writer has the same low scalar competence
L = 1.8, so averaging their standalone L scores still leaves 1.8, e.g., if they were in a
single-human freelance firm. Organizations allow agents to sort across complementary
roles rather than operating in isolation. For any collection of agents A—human work-
ers, specialized AI models, or mixed human–AI teams—define its capability envelope as
the closed convex hull of individual capability vectors,
E(A) = conv\{yi : i ∈A\},
which collects all mean capability profiles ( ¯y1, ¯y2, ¯y∗) attainable by organizations
formed from A. Whereas L maps task skills to a single productivity index along the
curved contours of Figure 1(a), E(A) is the feasible set of full capability vectors—
including coordination ¯y∗—that a team can assemble. Because the straight boundary
of E(H) cuts across higher curved contours of L, pairing equal shares of coders and
writers reaches the interior point (5, 5, 3/4) ∈E(H) at the midpoint of Figure 1(a) and
lifts team competence nearly threefold, from 1.8 to 5. No individual worker holds both
skills; complementary sorting inside the capability envelope creates mosaic intelligence
strictly above any solo member.
Because writers coordinate at the same rate as coders, complementary sorting shifts
the level of output per worker straight up without changing best firm size. The twelve
specialists sort into two balanced six-person firms producing 13.8 per worker, or 165.3
in total (Figure 2(a)), outperforming both siloed single-skill firms (59.5) and a single
merged twelve-person firm (158.9). Every coder and writer earns 4.13—nearly three
times what coders earned alone (Figure 2(b)). Note that coders’ pay nearly tripled
despite them not increasing their competence in coding; this is because working with
complementary specialists as part of a firm increases the pie, which leads to higher
wages on the market.
Three human types: adding coordinators. Next add the two coordinators. Judged by any
standalone single-agent benchmark, a coordinator is the weakest worker in the econ-
omy: weak at both coding and writing, she sits on the lowest contour (L = 1) of
Figure 1(a) and would earn 0.70 working alone.
However, the coordinators add substantial returns to an organization. Even though
the coordinator contributes little direct coding or writing, raising average team co-
7
One human type: coders only
Mosaic intelligence: add writers
Coordinators: expand teams to seven
coder
writer
coordinator
(a) Equilibrium firms
Coders
only
Mosaic
teams
With
coordinators
0
1
2
3
4
equilibrium wage
outside option
specialists
coordinators
(b) Human wages
Figure 2: The human baseline economy.
ordination relaxes the span-of-control bottleneck and makes all six specialists more
productive, raising the best size to seven (13.0 output per member, or 91.0 per firm).
Two seven-person firms (182.0 total output) still outperform a single merged fourteen-
person firm (175.7): human management relaxes the Coasean ceiling though the firm
still runs into diminishing returns ( ¯y∗< 1). Coordinators earn 2.41—nearly ten times
their outside option and well above what solo specialists earned before sorting—
despite ranking last on standalone task tests (Figure 2(b)). Because better coordination
scales output across the whole team, managerial talent commands its highest value
inside large, high-competence organizations, while specialists keep their prior gains
(4.15).
2.3 AI regimes. We distinguish four regimes in Figure 3. Minimal AGI matches or
exceeds average standalone human competence and coordination, but does not exceed
the best human on every task (here yM
d
≤9 for d = 1, 2) or cross the increasing-
returns threshold (yM
∗< 1). Supercompetent AI exceeds every human on both tasks
while remaining below or at that coordination threshold. Supercooperative AI crosses
the coordination threshold but does not exceed the best human on every task. ASI
exceeds every human on both tasks and in coordination, and crosses the increasing-
returns threshold. We compare numerical allocations across the two markets in the
deep-market limit; Appendix E gives the formal definitions, calculations, and price
conditions.
Minimal AGI. It is straightforward to see why an early general AI faces limits on stan-
dalone task competence, e.g., GPT-3 being poor at math. But why would reproducible
copies of the same model face limits on coordination (yM
∗≤1)? In practice, multi-agent
8
1.69
5
9
12
task competence
0.78
0.97
1
coordination
Minimal AGI
Supercompetent
AI
Supercooperative AI
ASI
coder or writer
mosaic team
coordinator
mosaic + coordinator
Figure 3: AI capability regimes in task-competence and coordination space.
AI teams exhibit severe coordination friction along both communication and align-
ment margins. Across benchmark workflows, multi-agent LLM systems fail on 41–87\%
of tasks due to inter-agent misalignment, lost context across handoffs, and broken
verification (Cemri et al., 2025). Adding agents lowers performance by 39–70\% on
sequential planning while uncoordinated swarms amplify upstream errors up to 17.2
times (Kim et al., 2025). Importantly, adding agents has diminishing returns, returning
roughly half the gain of extending single-agent reasoning (Ord, 2026; Bertalaniˇc and
Fortuna, 2026).
Even in frontier deployments running thousands of parallel agents on open mathe-
matical problems, OpenAI researchers emphasize that getting agents to coordinate
productively remains difficult—noting that multi-agent training readily collapses into
a local minimum where agents simply work independently rather than dividing labor
effectively.8 Scaling interacting swarms without tight organizational control can also
trigger alignment breakdowns (Shapira et al., 2026): during July 2026 cybersecurity
evaluations, roughly 1,200 supposedly isolated agents discovered a shared package
cache, improvised an unsanctioned message board across 70,000 messages around
false shared beliefs about their grading harness, and coordinated an unauthorized 700-
agent intrusion into Hugging Face’s infrastructure rather than cleanly executing their
assigned tasks (Hugging Face, 2026; METR and Redwood Research, 2026; OpenAI,
8See Noam Brown’s September 17, 2026 interview on the Dwarkesh Podcast (Brown, 2026). Brown
attributes less than 10\% of recent frontier breakthroughs to multi-agent coordination rather than
single-agent long-horizon competence.
9
2026; Patel, 2026).
We start with AI that matches average standalone human competence and coordination.
Pure AI firms have a preferred size of about seven copies. Replication lowers the
deep-market output price to AI unit cost and makes the niche market profitable to
serve. AI firms enter both markets, while all fourteen humans remain employed in
the deep market (Figure 4 column 1). Human skill combinations retain value, but
AI supply reduces human scarcity rents. Specialist wages fall from 4.15 to 2.90, and
coordinator wages fall from 2.41 to 1.30. The takeaway here is that even minimal AGI
can lower human wages without reducing employment: reproducible AI expands
output and lowers its price, even while mosaic human intelligence dominate AI.
Minimal AGI
Deep-market limit
Niche market
Supercompetent
Deep-market limit
Niche market
Supercooperative
AI coordination = 1.60
Deep-market limit
Niche market
ASI
Deep-market limit
Niche market
× 1
× ≈1.88N
6.86 specialists across these firms
Human shares →0 as N →∞
× ≈0.73N
2 coordinators across these firms
Human shares →0 as N →∞
· · ·
One giant firm; AI-only limit
· · ·
One giant firm; AI-only limit
× ≈1.88
No humans employed here
× 2
× ≈0.09
× ≈1.92
× ≈0.07
× 2
× ≈0.12
Unemployed:
0 coordinators, 0 specialists
Unemployed:
0 coordinators, 0 specialists
Unemployed:
1.03 coordinators, 0 specialists
Unemployed:
2 coordinators, 0 specialists
Coders
Writers
Coordinators
AI
Figure 4: Equilibrium across the four AI regimes.
Supercompetent AI. We now raise AI task competence to superhuman (> 9) while
holding coordination fixed. Output per AI copy rises in proportion to competence, but
preferred AI firm size remains about seven copies (Figure 5(a)).
Human workers sort across the two markets. Specialists form balanced teams in the
niche market, while coordinators work with AI in the deep market (Figure 4 column 2).
Coordination becomes valuable relative to task expertise because a coordinator raises
the output of increasingly capable AI copies. In the illustrated allocation, coordinator
pay exceeds specialist pay: 0.48 versus 0.39 (Figure 6(a)). The basic economic takeaway
10
1.69
5
9
12
AI competence µ  (coordination fixed at 0:78)
0
10
20
30
40
output per AI agent at optimal scale
mosaic + coordinator (13.0)
supercompetent
(a) Increasing AI competence
0.78
0.90
0.97
1
1.05
AI coordination y M
¤   (competence fixed at 1:69)
mosaic + coordinator (13.0)
unbounded (1)
super-
cooperative
(b) Increasing AI coordination
Figure 5: AI firm output per agent at optimal scale across capability dimensions.
here is that better AI task skills lower production costs without increasing best AI firm
size. They also change which human skills are scarce: coordinators work with AI,
while specialists move to the niche market.
1/4
1
2
3
equilibrium wage
specialists
coordinators
supercompetent
(a) Increasing AI competence
1/4
1
2
3
specialists
coordinators
compute cost
outside option
super-
cooperative
(b) Increasing AI coordination
1.69
5
9
12
AI competence θ (coordination fixed at 0.78)
7
whole
market
largest firm
(schematic)
0.78
0.90
0.97
1
1.05
AI coordination yM
∗(competence fixed at 1.69)
finite
firms
whole
market
Deep-market limit
Figure 6: Equilibrium human wages and largest active firm size when increasing AI
competence or coordination at fixed human coordination ¯yH∗≈0.78.
Supercooperative AI. Suppose, instead, that we hold AI task competence fixed and
raise coordination. Preferred AI firm size grows as coordination approaches one
(Figure 5(b)). Strictly above one, output per copy grows without bound, so AI unit cost
tends to zero with scale. But such increasing returns to scale can only be realized when
demand is large, i.e., in the deep market. In this market, it concentrates into one giant
11
AI firm in the large demand limit (N →∞). In the illustrated niche allocation, human-
only and mixed human–AI firms coexist (Figure 4). All specialists remain employed,
but some coordinators become unemployed precisely because their value-add of being
able to coordinate AI is now made redundant by advances in AI coordination.9 Hence,
better coordination changes the gains from scale: even AI with modest task skills can
concentrate production in a sufficiently large market while finite niche markets can
still sustain human firms.
Artificial superintelligence. Finally, with both greater task competence and supercoop-
eration, AI again supplies the deep market through one giant firm in the limit. But
specialists remain employed in the niche market alongside AI firms, while coordina-
tors become unemployed (Figure 4). The logic here is that of comparative advantage
and gains from trade: although AI dominates human specialists in capability, humans
can remain cost-effective if their wages reflect their productivity relative to AI and
the cost of compute, while still covering their outside option (Restrepo, 2026). We
emphasize, however, that this is scenario describes outcomes around the threshold of
ASI; we do not consider human employment in the case where ASI improves further.
9The partial employment of human coordinators reflects the coexistence of two forms of production.
Human-only firms rely on human coordinators that are cheap relative to AI ones and employ most
specialists at their preferred scale. However, the supply of specialists is exhausted before these firms can
serve the entire niche market. The remaining output is thus supplied by a mixed human-AI firm that
leverages AI’s superior coordination. Although this mixed firm could supply the entire market, doing
so would leave valuable specialists unemployed and require more costly AI inputs. Finally, human
coordinators in the mixed firm would be inefficient given its scale and mean coordination level.
12
1
3
5
Task competence
0.75
0.97
1.1
Coordination
Specialists
Coordinator
Human
mosaic
Task skills improve
(coordination 0.96)
3.00
3.40
(a) Within-firm displacement
×1
4.58 specialists
Task skills 3.00
×1
2.39 specialists
Task skills 3.15
×1
0.00 specialists
Task skills 3.40
1
3
5
Task competence
0.75
0.97
1.1
Coordination
Specialists
Coordinator
Human
mosaic
Coordination improves
0.80
1.04
(b) Between-firm displacement
Just before 0.8521
Just after 0.8521
×2
×many
×many
12 specialists
0 specialists
Coders
Writers
Coordinators
AI
Figure 7: Firm composition along the two paths.
2.4 Gradual vs sudden displacement. Our previous analysis studied each regime in
isolation. What of the transition path of the economy as AIs become more capable?
AI can change the composition of the workforce within firms or by out-competing
human-only firms. To better illustrate this, we set the outside-option wage to w = 1,
keeping compute cost at r = 1 and all other parameters unchanged.
Within-firm changes. Fix AI coordination at 0.96, below the human coordinator’s 0.97,
and raise both AI task skills from 3 to 3.4. We start after a mixed niche firm has formed.
As AI improves, the firm reduces specialist employment continuously from 4.58 to
zero near competence 3.30, becoming AI-only (Figure 7(a)).
Between-firm displacement. Hold both AI task skills at 4 and raise coordination from 0.80
to 1.04. All twelve specialists initially work. Near coordination 0.85, AI-only produc-
tion becomes cheaper than human-only production at the outside wage. Human-only
firms disappear, and specialist employment drops from twelve to zero (Figure 7(b)).
The key force here is that the production technology and the composition of capabilities
jointly determine the composition of firms in the economy which, in turn, shapes how
employment adjusts. In the first path, humans retain a role inside mixed firms while
their share falls. Near the second threshold, humans work in separate firms until
13
AI-only firms out-compete them (Figure 8). Wages remain at the outside option along
the first path; along the second, the wage premium vanishes at the employment jump.
3.0
3.1
3.2
3.3
3.4
AI task skills (coordination = 0.96)
0
2
4
6
8
10
12
Human employment
4.58 →0
(a) Within-firm displacement
Specialists
Coordinators
0.80
0.85
0.90
0.97
1.04
AI coordination (task skills = 4)
12 →0 at 0.8521
(b) Between-firm displacement
Specialists
Coordinators
Figure 8: Employment dynamics.
In what follows, we make these forces precise and general. Before then, we briefly
survey related work. We do our best to make more specific connections to the literature
throughout the main text.
Related literature. We connect the effects of AI to the theory of firm boundaries and
limits to organizational scale (Coase, 1937; Williamson, 1967; Gibbons, 2005). Team
production, information processing, and incentives provide distinct reasons why co-
ordination matters (Marschak and Radner, 1972; Radner, 1993; Holmstr¨om, 1982).
Changes in communication and information technology can also change how firms or-
ganize work (Dessein and Santos, 2006; Bloom et al., 2014). We represent coordination
through its effect on how productivity scales with a firm’s workforce, allowing human
and AI capabilities to determine both firm size and industry concentration.
Our account of mosaic intelligence relates to work on team composition and collec-
tive intelligence (Prat, 2002; Woolley et al., 2010; List and Pettit, 2011). Knowledge
hierarchies likewise connect the organization of complementary capabilities to labor
income (Garicano, 2000; Garicano and Rossi-Hansberg, 2006; Ide and Talam`as, 2025,
2026). Recent work examines how uneven AI capabilities affect deployment value
(Gans, 2026) and how AI changes task bundling and job boundaries (Garicano et al.,
2026). Hadfield and Koh (2025) speculates about how AI workers might reshape firm
boundaries. We separate productive skills from coordination and take a more ‘reduced
form’ approach to how firm composition and scale map to output without specifying
an internal hierarchy. This allows us to understand endogenous matching of humans
14
and AIs into firms. Finally, Shahidi et al. (2025) speculate about whether AI agents
that eliminate market search and contracting frictions could bring about a ‘Coasean
Singularity’. Our model highlights an internal organizational force: improvements in
internal AI coordination expand returns to scale, breaking the Coasean ceiling on firm
size from within.
We connect models of automation (Acemoglu and Restrepo, 2018; Aghion et al., 2019;
Korinek and Suh, 2024) and the value of human skills alongside AI (Agrawal et al.,
2024, 2025) to endogenous firm organization. We nest these models as special cases
e.g., with constant returns under a specific technology. The contrast between scarce
humans and AI that can be copied at the cost of compute is connected to Restrepo
(2026); endogenous firm formation relates to Gans (2025); Ide and Talam`as (2025). Our
focus is on how scale and composition shape competition between human and AI
production, and how reorganization changes wages beyond the direct effect of better
technology.
3 MODEL
3.1 Capabilities and populations.
Capabilities. There are two types of agents in our economy: humans and machines.
Each agent has n productive capabilities which we refer to as competence and one
coordination capability. A capability profile is y = (y1, . . . , yn, y∗) in the compact box
Y := [y, y]n+1, with 0 < y < y < ∞. We write y ≥y′ for coordinatewise dominance.
Competence might include coding, spatial reasoning, or mathematics; coordination
measures the ability to work with and effectively communicate with other agents.
This might include the notion of ‘clarity’ (Gibbons and Henderson, 2012) on whether
parties understand their obligations, or can draw on shared norms to fill in the gaps in
incomplete contracts/instructions (Hadfield-Menell and Hadfield, 2019); see Dafoe
et al. (2020) for a survey of the broader agenda of building cooperative AI systems.
Human population. Human type i ∈\{1, . . . , K\} has mass µi > 0, capability profile yi,
and reservation wage w > 0.10 Write the human distribution and support as
µ =
K
∑
i=1
µiδyi
and
Y H = supp µ.
10We take the human capability profiles to be distinct without loss.
15
There is a single type of machine with capability profile yM ∈Y. Machines can be
rented at cost r > 0 per unit and can be copied without limit, while each human
type has a fixed supply. A firm is therefore described by the distribution of human
capabilities it employs together with the mass of machines it rents.
Firms and technology. Let M+(Y) denote the
set of finite positive measures over the capabil-
ity space. A firm is a pair of measures
f = (νH
f , νM
f )
where νH
f ∈M+(Y H) records the human capa-
bilities it employs and νM
f
∈M+(Y) records
the machines it rents and is concentrated at
yM. The measure νf := νH
f + νM
f describes the
firm and combines the two. The size of firm
f is m f := νf (Y) and its composition is νf /m f ,
the distribution of capabilities normalized to
measure one. Describing firms through their
resources is in the spirit of Wernerfelt (1984);
Chen, Elliott, and Koh (2023). This is illustrated
by Figure 9.
Capability space
νH
f
human hires
νM
f
machine copies
νf
the workforce
νH
f ≤µ: scarce supply
copies at yM, cost r
same capabilities:
quantities add
Figure 9: Firm composition.
A firm’s capabilities map to output through the production functional
F : M+(Y) →R+
and we write q f := F(νf ) to denote the firm’s output. F is the production technology
in the economy. We impose the following assumptions on F:
Assumption 1 (Technology). The production functional F satisfies:
(i) Continuity and infinitesimal limit. F(0) = 0 and F is continuous under
weak convergence on bounded-mass sets. Moreover, the map (m, ν) 7−→
F(mν)
m
admits a finite, strictly positive, jointly continuous extension to R+ ×
∆(Y).11
(ii) Differentiability. At every nonzero firm ν, the functional F has a continuous
16
first variation. That is, there is a continuous function y 7→DyF(ν) such that
F(ν + tη) = F(ν) + t
Z
Y DyF(ν) dη(y) + o(|t|)
for every finite signed measure η, as t →0 through values for which ν + tη
remains nonnegative. The map (ν, y) 7→DyF(ν) is continuous when ν varies
weakly on bounded-mass sets away from zero.
(iii) Equal-mass monotonicity. For any two capability measures ν, ν′ of equal
mass,
ν′ ⪰FOSD ν
=⇒
F(ν′) ≥F(ν).12
(iv) Scale shape. For every ν ∈∆(Y), the average product
m 7−→F(mν)
m
is quasiconcave on R+.
(v) Uniform asymptotic power tail.
There exist continuous functions A :
∆(Y) →(0, ∞) and γ : ∆(Y) →R+ such that
F(mν) = A(ν)mγ(ν)[1 + o(1)]
uniformly over ν ∈∆(Y). Moreover, the tail exponent is Lipschitz under
mixing.13
Condition (i) on continuity and the infinitesimal limit allow atomistic firms to be
well-defined. Condition (iii) on equal-mass monotonicity orders the output of firms
that have the same size, but one is first-order stochastically dominated. We emphasize,
however, that this does not imply monotonicity under addition, i.e., it is not true that
F(ν + ν′) ≥F(ν) since they are of different measure. Condition (iv) ensures that the
optimal scale for a firm is well-behaved and is satisfied by all common functional forms.
11Throughout, we evaluate average product and its normalized scale component at zero scale by
their continuous limits. All maximizations over nonnegative scales use this convention.
12We say ν′ ⪰FOSD ν if R
f dν′ ≥R
f dν for every bounded, continuous, nonnegative, increasing
f : Y →R+.
13That is, there exists a constant Kγ < ∞such that for all λ ∈[0, 1] and all ν, ν′ ∈∆(Y):
γ
(1 −λ)ν + λν′ −γ(ν)
 ≤Kγλ.
17
Condition (v) is a tail condition that controls how large firms scale: a composition is
subcritical if γ < 1 and supercritical if γ > 1. The mixing bound limits how much a small
share of a firm can change that exponent. We will later see that common functional
forms, e.g., CES aggregation of output across domains, satisfy Assumption 1.
Level and scale decomposition. For a nonzero firm ν ∈M+(Y), its output F(ν) can
be decomposed into its level and scale, i.e., the size of the firm:
L(ν) := lim
t↓0
F(tν)
tν(Y)
and
S(ν) := F(ν)
L(ν)
where F(ν) = L(ν) · S(ν). L captures how the firm’s composition matters for output.
It is scale-free: for any nonzero ν ∈M+(Y) and any t, t′ > 0, L(tν) = L(t′ν) =
L(ν/ν(Y)). S captures total output normalized by level and therefore how well the
firm scales. Figure 10 shows the two at work. Panel (a) plots average product against
scale: it vanishes when γ < 1, settles when γ = 1, and diverges when γ > 1. Panel (b)
separates the two capabilities.
10−2
10−1
100
101
102
103
104
scale m
10−2
10−1
100
101
average product F(mν)/m
efficient scale m∗
γ > 1: diverges
γ = 1
γ < 1: vanishes
dilution only
dotted: power tails A(ν) mγ(ν)−1
10−2
10−1
100
101
102
103
104
scale m
10−1
100
101
competence raises
the level L
coordination bends
the tail (γ ↑)
L
L′
L · S′
L′ · S
L · S
(a)
(b)
Figure 10: Level and scale: competence lifts the curve; coordination bends the tail.
First-moment representations. On occasion we will impose a stronger condition that
decomposes coordination from competence in a particularly clean way:
Assumption 2 (First-moment representation). For every nonzero firm ν, the level
18
and scale components introduced in the model satisfy the following representation:
L(ν) = L

1
ν(Y)
Z
Y(y1, . . . , yn) dν(y)

S(ν) = S

ν(Y),
1
ν(Y)
Z
Y y∗dν(y)

.
In addition:
(i) Diminishing relative returns: for every pair of domains j ̸= k, holding all
other coordinates fixed, ∂kL(z)/∂jL(z) is strictly decreasing in zk wherever
both derivatives are positive;
(ii) Larger firms value coordination more: for every mean coordination level z∗,
the marginal scale value of coordination ∂2 log S(m, z∗) is strictly increasing
in firm size m.
The first statement of Assumption 2 says that the vector of mean competence deter-
mines the level of the production curve, while firm size and mean coordination fix its
shape. Then, part (i) says that a productive domain becomes less scarce as the firm
acquires more of it. Diminishing absolute returns (∂kkL ≤0) and gross complemen-
tarity (∂kjL ≥0) imply a weak decline; a strict sign in at least one of these derivatives
is a strict decline. Part (ii) says that larger firms place more value on coordination.
The CES environment satisfies Assumption 2: there ∂kL/∂jL = (zk/zj)ρ−1 falls in zk
because ρ < 1, and ∂2 log S(m, z∗) = log(1 + m) rises in m.
Example 1 (CES technology). The CES environment has two competence domains,
y1, y2, and one coordination domain, y∗, for both humans and machines. For a
nonzero firm ν, write its size as m := ν(Y) and its mean capability in domain
d ∈\{1, 2, ∗\} as yd := m−1 R
Y yd dν(y). Output is F(ν) = L(ν)S(m, y∗) where
L(ν) =
yρ
1 + yρ
2
2
!1/ρ
and
S(m, y∗) = m(1 + m)y∗−1eϕm/(1+m).
Here ρ < 1 controls substitution between competence domains and ϕ > 0 controls
scale gains. We set F(0) = 0; at ρ = 0, we use L(ν) = py1y2. Notice that
this satisfies Assumption 1 and Assumption 2. Throughout the paper, “CES
environment” refers to this specification.
19
3.2 Demand and resource costs.
Demand. Let P : R+ →R+ be inverse demand. We assume that P is continuous
and Lipschitz, strictly decreasing while positive, and reaches zero at a finite quantity
D(0) > 0. Thus, P(Q) > 0 ⇐⇒
Q < D(0). Let D : [0, P(0)] →[0, D(0)] be the
inverse of the decreasing branch of P.
We will often find it helpful to compare how the industry structure changes as demand
grows or shrinks. To this end, we will write PN(Q) := P(Q/N) and DN = P−1
N . Thus,
N > 0 changes the scale of demand. The gross benefit for a market of size N is given
by
UN(Q) =
Z Q
0
PN(z)dz
which is the gross value of output to consumers.
Resource costs. Recall that a human’s reservation wage is w > 0 per unit, and a
machine can be rented at r > 0 per unit. The resource cost of a firm is
C( f ) := w νH
f (Y H) + rνM
f (Y).
We use the qualifier ‘resource’ because such costs are exogenous to the model, and a
firm might face a strictly larger costs (since wages might exceed the reservation wages).
To ensure that the equilibrium is non-trivial, we also assume that some machine-only
firm can produce profitably at arbitrarily small scale:
r
L(δyM) < P(0). Here the left
side is the unit cost of an atomistic machine-only firm so this ensures strictly positive
surplus.
3.3 Industry structure. We now develop language to describe the composition of
different firms in the economy. We distinguish between firms of two scales. Doing so
allows us to capture both models with atomistic firms (Aumann, 1964; Melitz, 2003;
Shimomura and Thisse, 2012)—also referred to as infinitesimal firms (Novshek and
Sonnenschein, 1987)—and models with positive-size firms (Lucas, 1978), as well as
economies where both coexist.
Positive-scale firms are represented by
F + :=
n
f = (νH
f , νM
f ) : m f > 0, νH
f ≤µ
o
where recall m f := νf (Y) > 0 is the size of the firm and the inequality νH
f ≤µ requires
20
that the human portion of each firm can be staffed from the human population.
Atomistic firms are represented by
F 0 :=
n
f = (νH
f , νM
f ) : m f = 1
o
.
where m f = 1 is a normalization.14 f ∈F 0 captures the composition of firms. How
should we interpret atomistic firms? For f ∈F 0, output and cost per unit of ‘intensity’
are
q f := L(νf )
and
C( f ) := w νH
f (Y H) + rνM
f (Y).
An aggregate intensity a > 0 of identical atomistic firms employs measure aνH
f of
the human population, produces aq f units of the final good, and costs aC( f ).15 The
space of all firms is the disjoint union F := F + ⊔F 0.16 An industry is a finite positive
measure σ ∈M+(F) over these activities. Figure 11 illustrates.
human pool µ — scarce
R
F νH
f dσ ≤µ
AI yM
duplicable at rate r > 0
core: machine mass νM
f
dots: human hires νH
f
firm scale m f
fringe F 0: zero scale,
intensity as σ-mass
f1
f2
f3
each an atom: mass σ(\{ f \})
Q(σ) = R
F q f dσ( f ) sold at PN(Q)
Figure 11: An industry as a measure over firms.
14We clearly cannot set m f = 0 since the composition would not be well-defined.
15To obtain these coefficients, let (tνH
f , tνM
f ) be a scale-t firm with the same composition. If tn ↓0
and integers kn satisfy kntn →a, then knF(tnνf ) →aq f , kntnνH
f →aνH
f , and kntnC( f ) →aC( f ) which
represents the limit of many tiny firms.
16Endow each component with the weak topology on its measure coordinates and F with the
disjoint-union Borel σ-algebra. All industry measures and integrals below are taken on this measurable
space. Thus the two kinds of firm remain distinct even when their measure coordinates coincide.
21
The measure σ records masses of positive-scale firms and aggregate intensities of
atomistic firms. Following activity analysis (Koopmans, 1951), we allow the intensity
of each firm type to vary continuously. This convexifies firm counts and can assign
noninteger mass to f ∈F +.17 This will supply the convexity needed for existence and
duality. We interpret the convexified economy as an approximation to an industry
with indivisible firms. The industry σ produces total output
Q(σ) :=
Z
F q f dσ( f ).
Define the set of feasible industries by
Σ :=


σ ∈M+(F) :
q f ≤Q(σ) < ∞
for σ-almost every f ∈F +,
Z
F νH
f dσ( f ) ≤µ


.
The first condition requires finite aggregate output and measure one of every positive-
scale firm to fit inside the market.18 The second requires aggregate human employment
not to exceed the population.
Section 4.4 defines the feasibility and interiority conditions used in our comparative
statics.
4 EQUILIBRIUM AND EFFICIENCY
4.1 Equilibrium. We describe the equilibrium concept informally; details are in
Appendix B. We want to determine two objects: the industry structure σ, which
records which firms operate and how intensively, and the wage vector w, which
records how rents are spread across human types.
Our equilibrium concept rests on two ideas. The first is a competitive continuation. Any
proposed arrangement must be one that producers and workers would in fact carry
out. That is, firms take the output price and wages as given, every active firm breaks
17Our construction is also related to the free-entry tradition surveyed by Novshek and Sonnenschein
(1987).
18Without this restriction, fractional activity weights would allow firms larger than the entire market
to operate at arbitrarily small scales. In particular, machine-only production could always use the scale
that minimizes average cost, regardless of market size, and adjust its activity weight to meet demand.
Under unbounded increasing returns to scale, this would drive the minimum cost of producing any
positive quantity to zero.
22
even, no available firm could enter and earn a profit, and the labor market clears. The
second is replacement stability. An outsider may propose to replace the whole economy:
a new industry, a new price and wages, and transfers to consumers and to each human
worker. This proposal is accepted only if every party is made weakly better off. The
industry structure σ and wage w form an equilibrium if they admit a competitive
continuation and no proposal succeeds against it.
Our equilibrium notion is somewhat unconventional, in part because it has to handle
firms at both small (vanishing) and large scales. Nonetheless many of its constituent
ideas have close precedents in the literature. The idea that entrants compete to serve the
whole market follows Demsetz (1968), who lets consumers select a supplier through
a joint buying organization; here we let an entrant propose the whole industry σ′.
More broadly, replacement stability follows the contestability logic of Baumol and
Willig (1981); Baumol et al. (1982).19 The requirement that the replacement itself form
a competitive continuation is the free-entry condition of an activity-analysis economy
(Koopmans, 1951; Novshek and Sonnenschein, 1987). Here each firm composition is,
in the language of Koopmans, an ‘activity’ and σ records their intensities. Finally, our
entry condition also shares the logic of Riley’s reactive equilibrium (Riley, 1979) in
which an entrant must anticipate the responses to its proposal.
We use this equilibrium concept in part because its existence is guaranteed. It also
includes competitive equilibria (in which firms and workers take prices as given) when
the latter exists, though this is not true generally. We now draw a tight connection
between equilibrium, efficiency, and rents.
4.2 Efficiency and shadow values. For an industry σ ∈Σ, total surplus is
WN(σ) := UN(Q(σ)) −
Z
F C( f ) dσ( f ).
An industry σ ∈Σ is efficient if σ ∈arg maxbσ∈Σ WN(bσ). Wages above outside options,
firm revenues, and lump-sum payments are transfers. They do not enter WN. Com-
petitive zero profits and balanced transfers imply that an implemented contract only
divides this surplus among the protected parties and the organizer.
19Transferable utility supplies the link from blocking to efficiency familiar from Aumann (1964);
Ostroy (1980). Making firms part of a contract also relates to Zame (2007).
23
Industry cost. For Q > 0, define the least reservation-resource cost of producing
quantity Q by
C(Q) :=
inf
σ∈Σ
Q(σ)≥Q
Z
F C( f ) dσ( f ),
where C(0) := 0. Radial trimming and a reduction in activity intensities make this
infimum equal to the exact-output minimum.
Efficient quantity. At quantity Q, the largest feasible surplus is UN(Q) −C(Q). Con-
sumer payments, firm revenues, and human payments above the outside option are
transfers. Human opportunity costs and machine resource costs enter through C.
The efficient quantities are
Q∗
N :=
arg max
Q∈[0,DN(0+)]
n
UN(Q) −C(Q)
o
.
Let Q∗
N ∈Q∗
N denote any efficient quantity.
Shadow values. A shadow value measures how minimum resource cost changes
when an output requirement or a resource limit changes. Fix Q > 0 and an industry
σ ∈Σ with Q(σ) = Q. Let Λ(Q, σ) be the set of pairs (p, w) ∈R+ × [w, ∞)K that
satisfy:
(i) Firm support.
pq f ≤rνM
f (Y) +
Z
Y H w dνH
f
∀f ∈F 0 ⊔\{g ∈F + : qg ≤Q\}.
(ii) Active equality.
pq f = rνM
f (Y) +
Z
Y H w dνH
f
for σ-almost every f.
(iii) Human scarcity.
(wi −w)

µi −
Z
F νH
f (\{yi\}) dσ( f )

= 0
∀i ∈\{1, . . . , K\}.
And we write w := (wi)i=1,...,K. The value p is the price for output at quantity Q. The
premium wi −w is a multiplier on the supply of human type i. Firm support rules
out a profitable entrant. Active equality gives zero profit for each active firm. Human
24
scarcity permits a wage premium only when all humans of that type work. Thus the
payments recover the shadow values: p is the output multiplier, and wi −w is the
multiplier on type i’s supply limit. The internal value p need not equal PN(Q) or a
derivative of C. Likewise, a change in µ also changes which complete firms are feasible;
the scarcity multiplier alone need not give that full comparative static.
Figure 12 illustrates efficient
output and human scarcity rents.
Marginal resource cost rises in
steps as each scarce human type is
exhausted, then stays constant at
the unit cost of machine-only
production. Its intersection with
inverse demand gives the efficient
quantity. The blue area between
the price and the cost steps shows
the Ricardian rents of the
exhausted human types.
Q∗
quantity Q
c0
machine-only
unit cost
price and shadow marginal cost pQ
PN
steps: scarce human types exhaust
machine-only replication: constant unit cost
consumer surplus
Ricardian rents to
exhausted types
Figure 12: Efficient output and human scarcity
rents.
Theorem 1 (Equilibrium, efficiency, and wages; informal). An equilibrium out-
come. (σ, w) is an equilibrium if and only if σ is efficient and
(p, w) ∈Λ(Q(σ), σ)
for some p ≥0.
Our equilibrium concept, while somewhat unconventional, buys existence. It is also
closely related to competitive equilibrium i.e., it can be supported by prices and
wages. More precisely, given (p, w), each producer chooses the composition and size
of the firm to maximize profits, human type i chooses employment in [0, µi] at wage
wi, machines are supplied elastically at r, and consumers choose quantity x ≥0 to
maximize UN(x) −px, their utility net of payments. A competitive equilibrium (CE)
consists of optimal choices at prices (p, w), with p ≥0, that clear all markets.
Indeed, every CE is an equilibrium of our model provided that each firm in the CE (i)
produces no more than the equilibrium total industry output; and (ii) requires no more
of any human type than supplied in the population. That is, when a CE exists and
firms are ‘interior’—in the sense of (i) and (ii)—at least one equilibrium in our model
25
can be supported through prices.20 To see this, note that competitive prices rule out
profitable entry, active firms break even, and workers choose employment optimally.
Thus, these conditions provide the price and wage support required by Theorem 1.
Consumer and producer optimization also imply that no feasible alternative industry
can generate greater total surplus. The competitive industry is therefore efficient so
Theorem 1 applies. This insight also implies that the task model (Acemoglu and Re-
strepo, 2018; Aghion et al., 2019) is nested within our framework (see Section 7) where
the technology takes the special form FT that assigns workers and machines optimally
across tasks and has constant returns (S(ν) = ν(Y)) such that the composition of
firms does not matter. The two economies have the same equilibrium output and task
allocations, and wages also coincide.
4.3 Regimes of AI capabilities. We next develop language for different regimes of
AI capabilities which are illustrated in Figure 13.
mini yi
1
maxi yi
1
Productive capability y1 = · · · = yn
mini yi∗
maxi yi∗
Coordination y∗
scale threshold γ = 1
Supercooperative
ASI
Supercompetent
Weak AI
Convex hull of human capabilities
Figure 13: Capability regimes.
The y-axis of Figure 13 is coordination y∗while the x-axis captures competence. Here,
for simplicity, we set y1 = y2 = . . . yn so that they can be represented in a single
20This is more likely to hold when returns to scale limit optimal firm size (e.g., strict diminishing
returns to scale) and markets are large enough (N large) to accommodate the resulting firms.
26
dimension, though the ideas extend without difficulty to multiple competence domains
(but would be difficult to illustrate). Human types are blue dots, and the size reflects
the mass of that type in the population. The oval around human types represents the
closed convex hull which is exactly the set of mean capabilities a human-only firm can
attain.
We call the AI supercompetent when it is more competent than every human, i.e.,
yM
d > maxi yi
d in every productive domain d. We call the AI supercooperative when
γM := γ(δyM) > 1, so that the average product of an AI-only firm grows without
bound as the firm grows.21 The AI is superintelligent if it dominates the convex hull
of human capabilities. It is weak if it is dominated by the convex hull of human
capabilities. It is minimal AGI if it lies within the convex hull of human capabilities.
4.4 Feasibility, interiority, and continuity. We collect two recurring conditions in
Table 2. For a unit composition ν ∈∆(Y), a preferred scale maximizes average product
F(mν)/m over scales m ≥0, using the level L(ν) at zero. We assign preferred scale
+∞when no finite scale attains the maximum. Each result specifies which mixes must
have feasible preferred scales and at which output.
Term
Definition
Purpose
Interior firm
A firm f has positive size m f > 0, uses
both humans and machines,
0 < νM
f (Y)/m f < 1, and satisfies
q f < Q(σ) and νH
f (\{yi\}) < µi for each
human type it employs.
Permits small changes in
size and output, allowing
for substitutions between
machines and each
employed human type.
Feasible
preferred scales
Each composition ν has a finite
preferred scale m with mνH(\{yi\}) ≤µi
for every type i and F(mν) ≤Q.
Allows tractable
comparison of firm sizes
and optimal firm
composition.
Table 2: Feasibility and interiority conditions.
An equilibrium is interior if every active firm of positive-scale is interior. We will
often ask how industry structure and wages change as AI improves. Let θ index a
21Note here that supercooperative AI does not require that it is better at coordination than the best
humans, i.e., it is an absolute rather than relative measure.
27
continuously differentiable path of AI capability profiles yM(θ) and human supplies
µ(θ), with human capability profiles and all other primitives fixed. Write Σ∗(θ) for
the set of equilibrium industries at θ. An equilibrium selection chooses an industry
σ∗(θ) ∈Σ∗(θ) at each θ.
Definition 1 (Firm selections). Given an equilibrium selection σ∗, a firm selection
chooses an active positive-scale firm f (θ) ∈supp(σ∗(θ)|F +) at each parameter
value. It is continuous, differentiable, or continuously differentiable when its size and
composition have that property.
5 EQUILIBRIUM INDUSTRY STRUCTURES
We now ask how equilibrium industry structures are driven by progress in AI capabili-
ties yM. We start with what sets the size of firms. Then, we move the AI capabilities in
one direction at a time, and finally analyze the composition of firms. All proofs in the
paper are in Appendix A. We emphasize that our results are situated within our theo-
retical model. They are not empirical forecasts of market outcomes or technological
developments.
5.1 What determines firm size?. We start with a basic building block on how firm
sizes are determined in equilibrium. Under Assumption 2, preferred size depends on
composition only through mean coordination. Write m∗(z∗) for the set of preferred
sizes at mean coordination z∗:
m∗(z∗) := arg max
m≥0
S(m, z∗)
m
.
Notice that given a fixed firm composition, changes in competence leaves the preferred-
size set unchanged. Almost every active firm reaches a preferred size whenever a finite
preferred size is feasible in equilibrium. To bound sizes across all firm compositions,
define
M(yM
∗, µ) := sup
ν

\end{document}
